\documentclass[main.tex]{subfiles}


\begin{document}

%from pagenumber to up
% \phantomsection 
% \hypertarget{toc}{}

 

 

% номер секции вручную
\setcounter{section}{5
}
\addtocounter{section}{-1} 

\section{Линза}


\textbf{Линза}~--- это прозрачное тело, ограниченное сферическими поверхностями. Выделяют два вида линз\footnote{Далее предполагается, что показатель преломления линзы \emph{больше} показателя преломления окружающей среды.}.
\begin{enumerate}
    \item \textbf{Собирающая}: толщина середины линзы \emph{больше} толщины ее края. На рис.\,\ref{fig:p180} показаны все возможные собирающие линзы.

    \begin{figure}[H]
    \centering

    \begin{tikzpicture}[font=\small,            line cap=round,                             >={Stealth[inset=0pt,round,length=3mm, angle'=20]},vec/.style={>={Stealth[inset=0pt,round,length=3.5mm, angle'=20]}}]


\pgfmathsetmacro{\lensRadius}{1.5}
\pgfmathsetmacro{\lensHeight}{1}
\pgfmathsetmacro{\startAngle}{asin(\lensHeight/\lensRadius)}

\fill[SkyBlue,semitransparent] (0,\lensHeight) arc[start angle=180-\startAngle,delta angle=2*\startAngle,radius=\lensRadius]
arc[start angle=-\startAngle,delta angle=2*\startAngle,radius=\lensRadius];


\draw (0,\lensHeight) arc[start angle=180-\startAngle,delta angle=2*\startAngle,radius=\lensRadius] node[below] {Двояковыпуклая}
arc[start angle=-\startAngle,delta angle=2*\startAngle,radius=\lensRadius];



%%%


\begin{scope}[xshift=4.5cm]

\pgfmathsetmacro{\lensRadius}{1.5}
\pgfmathsetmacro{\lensHeight}{1}
\pgfmathsetmacro{\startAngle}{asin(\lensHeight/\lensRadius)}

\fill[SkyBlue,semitransparent] (0,\lensHeight) -- (0,-\lensHeight)
(0,\lensHeight) arc[start angle=\startAngle,delta angle=-2*\startAngle,radius=\lensRadius];


\draw (0,\lensHeight) -- (0,-\lensHeight) node[below] {Плосковыпуклая};
\draw (0,\lensHeight) arc[start angle=\startAngle,delta angle=-2*\startAngle,radius=\lensRadius];

\end{scope}


%%%


\begin{scope}[xshift=9cm]

\pgfmathsetmacro{\lensRadius}{1.5}
\pgfmathsetmacro{\lensHeight}{1}
\pgfmathsetmacro{\startAngle}{asin(\lensHeight/\lensRadius)}

\pgfmathsetmacro{\lensRadiusb}{3}
\pgfmathsetmacro{\lensHeightb}{1}
\pgfmathsetmacro{\startAngleb}{asin(\lensHeightb/\lensRadiusb)}

\fill[SkyBlue,semitransparent] (0,\lensHeight) arc[start angle=\startAngle,delta angle=-2*\startAngle,radius=\lensRadius]
arc[start angle=-\startAngleb,delta angle=2*\startAngleb,radius=\lensRadiusb];


\draw (0,\lensHeight) arc[start angle=\startAngle,delta angle=-2*\startAngle,radius=\lensRadius] node[below,xshift=.3cm] {Вогнуто"=выпуклая}
arc[start angle=-\startAngleb,delta angle=2*\startAngleb,radius=\lensRadiusb];

\end{scope}

    \end{tikzpicture}
\caption{Собирающие линзы}
\label{fig:p180}
\end{figure}
\nointerlineskip

    \item \textbf{Рассеивающая}: толщина середины линзы \emph{меньше} толщины ее края. На рис.\,\ref{fig:p181} изображены все возможные рассеивающие линзы.

    \begin{figure}[H]
    \centering

    \begin{tikzpicture}[font=\small,            line cap=round,                             >={Stealth[inset=0pt,round,length=3mm, angle'=20]},vec/.style={>={Stealth[inset=0pt,round,length=3.5mm, angle'=20]}}]


\pgfmathsetmacro{\lensRadius}{2}
\pgfmathsetmacro{\lensHeight}{1}
\pgfmathsetmacro{\startAngle}{asin(\lensHeight/\lensRadius)}

\fill[SkyBlue,semitransparent] (0,\lensHeight) --++ (.4,0) arc[start angle=180-\startAngle,delta angle=2*\startAngle,radius=\lensRadius]
--++ (-.8,0)
arc[start angle=-\startAngle,delta angle=2*\startAngle,radius=\lensRadius]
--++ (.4,0)
;


\draw (0,\lensHeight) --++ (.4,0) arc[start angle=180-\startAngle,delta angle=2*\startAngle,radius=\lensRadius]
--++ (-.4,0) node[below] {Двояковогнутая}
--++ (-.4,0)
arc[start angle=-\startAngle,delta angle=2*\startAngle,radius=\lensRadius]
--++ (.4,0)
% --++ (0,.4) --++ (0,-.4)
;

%%%


\begin{scope}[xshift=4.5cm]

\pgfmathsetmacro{\lensRadius}{2}
\pgfmathsetmacro{\lensHeight}{1}
\pgfmathsetmacro{\startAngle}{asin(\lensHeight/\lensRadius)}

\fill[SkyBlue,semitransparent] (0,\lensHeight) -- (0,-\lensHeight)
--++ (.6,0)
arc[start angle=-180+\startAngle,delta angle=-2*\startAngle,radius=\lensRadius]
--++ (-.6,0)
;


\draw (0,\lensHeight) -- (0,-\lensHeight)
--++ (.3,0) node[below] {Плосковогнутая}
--++ (.3,0)
arc[start angle=-180+\startAngle,delta angle=-2*\startAngle,radius=\lensRadius]
--++ (-.6,0)
;

\end{scope}


%%%


\begin{scope}[xshift=9cm]

\pgfmathsetmacro{\lensRadius}{3}
\pgfmathsetmacro{\lensHeight}{1}
\pgfmathsetmacro{\startAngle}{asin(\lensHeight/\lensRadius)}

\pgfmathsetmacro{\lensRadiusb}{1.3}
\pgfmathsetmacro{\lensHeightb}{1}
\pgfmathsetmacro{\startAngleb}{asin(\lensHeightb/\lensRadiusb)}

\fill[SkyBlue,semitransparent] (0,\lensHeight) arc[start angle=180-\startAngle,delta angle=2*\startAngle,radius=\lensRadius]
--++ (.6,0)
arc[start angle=-180+\startAngleb,delta angle=-2*\startAngleb,radius=\lensRadiusb]
--++ (-.6,0)
;


\draw (0,\lensHeight) arc[start angle=180-\startAngle,delta angle=2*\startAngle,radius=\lensRadius]
--++ (.3,0) node[below,xshift=-.1cm] {Выпукло"=вогнутая}
--++ (.3,0)
arc[start angle=-180+\startAngleb,delta angle=-2*\startAngleb,radius=\lensRadiusb]
--++ (-.6,0)
;

\end{scope}



    \end{tikzpicture}
\caption{Рассеивающие линзы}
\label{fig:p181}
\end{figure}

\end{enumerate}
\nointerlineskip

% На рис.\,\ref{fig:p182} показаны качественные изменения пучка света при прохождении собирающей и рассеивающей линз. 

Пусть имеются две линзы; на них падают параллельные пучки света (рис.\,\ref{fig:p182}).

\tikzfading[name=fade right,
left color=transparent!0,
right color=transparent!90]

\tikzfading[name=fade rightb,
left color=transparent!40,
right color=transparent!70]

    \begin{figure}[H]
    \centering

    \begin{tikzpicture}[font=\small,            line cap=round,                             >={Stealth[inset=0pt,round,length=3mm, angle'=20]},vec/.style={>={Stealth[inset=0pt,round,length=3.5mm, angle'=20]}}]


\pgfmathsetmacro{\lensRadius}{3}
\pgfmathsetmacro{\lensHeight}{1}
\pgfmathsetmacro{\startAngle}{asin(\lensHeight/\lensRadius)}


%light
\fill[yellow,path fading=fade right] (-2,.4) -- (0,.4) -- (2,0) -- (0,-.4) -- (-2,-.4) -- cycle;
\draw[orange,->] (-2,.6) --++ (1,0) node[midway,above,black] {Свет};


\fill[SkyBlue,semitransparent] (0,\lensHeight) arc[start angle=180-\startAngle,delta angle=2*\startAngle,radius=\lensRadius]
arc[start angle=-\startAngle,delta angle=2*\startAngle,radius=\lensRadius];


\draw (0,\lensHeight) arc[start angle=180-\startAngle,delta angle=2*\startAngle,radius=\lensRadius]
arc[start angle=-\startAngle,delta angle=2*\startAngle,radius=\lensRadius];


%%%


\begin{scope}[xshift=6cm]

\pgfmathsetmacro{\lensRadius}{2.5}
\pgfmathsetmacro{\lensHeight}{1}
\pgfmathsetmacro{\startAngle}{asin(\lensHeight/\lensRadius)}


%light
\fill[yellow,path fading=fade right] (-2,.4) -- (0,.4) -- (2,.8) -- (2,-.8) -- (0,-.4) -- (-2,-.4) -- cycle;
\draw[orange,->] (-2,.6) --++ (1,0) node[midway,above,black] {Свет};


\fill[SkyBlue,semitransparent] (0,\lensHeight) --++ (.3,0) arc[start angle=180-\startAngle,delta angle=2*\startAngle,radius=\lensRadius]
--++ (-.6,0)
arc[start angle=-\startAngle,delta angle=2*\startAngle,radius=\lensRadius]
--++ (.3,0)
;


\draw (0,\lensHeight) --++ (.3,0) arc[start angle=180-\startAngle,delta angle=2*\startAngle,radius=\lensRadius]
--++ (-.3,0)
--++ (-.3,0)
arc[start angle=-\startAngle,delta angle=2*\startAngle,radius=\lensRadius]
--++ (.3,0)
;




\end{scope}

    \end{tikzpicture}
\caption{Действия собирающей и рассеивающей линз}
\label{fig:p182}
\end{figure}


Собирающая линза преобразует параллельный пучок света в \emph{сходящийся} пучок (рис.\,\ref{fig:p182}, слева), а рассеивающая~--- в \emph{расходящийся} (рис.\,\ref{fig:p182}, справа).

Линзу можно считать \emph{тонкой}, если \emph{ее толщина много меньше характерных расстояний в задаче.} На рис.\,\ref{fig:p183} показаны обозначения тонких линз. 

    \begin{figure}[H]
    \centering

    \begin{tikzpicture}[use optics,font=\small,            line cap=round,                             >={Stealth[inset=0pt,round,length=3mm, angle'=20]},vec/.style={>={Stealth[inset=0pt,round,length=3.5mm, angle'=20]}}]


\node[lens,thick,line cap=butt] (L1) at (0,0) {};
\node[below] at (L1.south) {Собирающая};

\node[lens,thick,lens type=diverging,line cap=butt] (L1) at (4,0) {};
\node[below,yshift=-4pt] at (L1.south) {Рассеивающая};


    \end{tikzpicture}
\caption{Обозначения тонких собирающей и рассеивающей линз}
\label{fig:p183}
\end{figure}















 
\end{document}